\chapter{1920:Walther Nernst}
\label{chap:chemistry1920}

The Nobel Prize in Chemistry for the year 1920 was awarded to Walther Nernst.

\textbf{Motivation:}
in recognition of his work in thermochemistry

\section{Walther Nernst}

\subsection*{Early Life and Education}

Walther Nernst was born on 1864-06-25 in Briesen, Prussia (now Wąbrzeźno, Poland).
Walther Hermann Nernst was a German physical chemist known for his work in thermodynamics, physical chemistry, electrochemistry, and solid-state physics. His formulation of the Nernst heat theorem helped pave the way for the third law of thermodynamics, for which he won the 1920 Nobel Prize in Chemistry. He is also known for developing the Nernst equation in 1887.

\subsection*{Career and Major Research}

Nernst was professor of physical chemistry at Göttingen and later in Berlin, where he directed the Physikalisch-Chemisches Institut. He formulated the heat theorem, stated in 1906, which provides a basis for calculating chemical equilibria at low temperatures and became known as the third law of thermodynamics. He also derived the equation for the electrode potential of an electrochemical cell that bears his name.

\subsection*{Nobel Prize-winning Work}

The work for which Walther Nernst was awarded the Nobel Prize in Chemistry in 1920 was described by the Nobel Committee as: "in recognition of his work in thermochemistry".

This research fundamentally advanced the field by making groundbreaking contributions that transformed chemical science.

\subsection*{Significance and Impact}

The impact of Walther Nernst's work extends far beyond the specific achievement recognized by the Nobel Prize.
These contributions continue to influence chemical research and education today.

\subsection*{Later Career and Honors}

Walther Nernst continued his scientific work until 1941-11-18, passing away in Zibelle, Germany (now Niwica, Poland).
He received numerous honors including membership in major scientific academies and societies.

\subsection*{Legacy}

The legacy of Walther Nernst endures through the lasting impact of his scientific contributions.
Walther Nernst's work continues to be cited and built upon by researchers across the chemical sciences.

\subsection*{Selected Publications}

Walther Nernst's major publications include numerous papers in leading journals of the era, detailing the experimental and theoretical work summarized above.

\clearpage

\section*{Chapter Summary}

Walther Nernst received the prize in recognition of his work in thermochemistry, including the heat theorem that bears his name.
